\numpara{3.4.7}\label{IV.3.4.7}%
{\renewcommand{\thefootnote}{(16)}\nde{: One has added this paragraph.}}%
Suppose that \(\mathcal{C}\) possesses a final object \(e\) and let
\(f\colon G\to G'\) be a morphism of \(\mathcal{C}\)-groups, such that
\(f\in(M)\). Then, by~3.4.1\,(a), the kernel \(\Ker(f)\) is representable
by \(e\times_{G'}G\), and the morphism \(\Ker(f)\to e\) belongs to~\((M)\).

On the other hand, the equivalence relation defined by~\(f\) is the same as
that defined by the action of \(\Ker(f)\) (say, on the right) on~\(G\),
\ie it is the image of the morphism
\(G\times\Ker(f)\to G\times G\), defined set-theoretically by
\((g,h)\mapsto(g,gh)\). Consequently, one deduces from~3.3.2.1 the
following corollary:

\begin{corollary}{3.4.7.1}
\label{IV.3.4.7.1}
Suppose that \(\mathcal{C}\) possesses a final object \(e\) and let
\(f\colon G\to G'\) be a morphism of \(\mathcal{C}\)-groups, such that
\(f\in(M)\). Then the action of \(\Ker(f)\) on \(G\) is \((M)\)-effective
and \(G'\) is the quotient \(G/\Ker(f)\).
\end{corollary}

\par\addvspace{0.55\baselineskip}
\noindent\textbf{3.5. Construction of quotients by descent.}\label{IV.sec.3.5}\ ---
It frequently happens that one does not know how to construct a quotient
directly, but that one knows how to do so after a suitable base change.
The present number gives a useful criterion in this situation.

One has seen in paragraph~2.1 the definition of a descent datum on an
object \(X'\) over \(S'\) relative to a morphism \(S'\to S\).

\begin{definition}{3.5.1}
\label{IV.3.5.1}
One says that a descent datum on \(X'\) relative to \(S'\to S\) is
\emph{effective}, if \(X'\) equipped with this descent datum is isomorphic
to the inverse image on \(S'\) of an object \(X\) over \(S\), equipped with
its canonical descent datum.
\end{definition}

If \(S'\to S\) is a descent morphism (2.2), then the \(X\) of the definition
is unique up to unique isomorphism. To say that \(S'\to S\) is an effective
descent morphism (2.2) is to say that it is a descent morphism and that
every descent datum relative to this morphism is effective.

Let us now consider an equivalence relation \(R\) in an object \(X\) over
\(S\). Let \(X'\) (\resp \(X''\), \resp \(X'''\)) be the inverse images of
\(X\) on \(S'\), \(S''=S'\times_{S}S'\) and
\(S'''=S'\times_{S}S'\times_{S}S'\), and let \(R'\), \(R''\), \(R'''\) be
the equivalence relations deduced from \(R\) by inverse image. Suppose
that the equivalence relation \(R'\) in \(X'\) is \((M)\)-effective, and
consider the quotient \(Y'=X'/R'\), which is an object over \(S'\). Its
two inverse images on \(S''\) are isomorphic to \(X''/R''\) by~3.4.5. The
\(S'\)-object \(Y'\) is therefore equipped with a canonical gluing datum.
Using likewise the uniqueness of \(X'''/R'''\), one sees that it is even a
descent datum. (Remark: one has implicitly supposed in this proof that
all the fiber products written existed, which is the case in particular if
\(S'\to S\) is quarrable, for example a descent morphism.)

\begin{proposition}{3.5.2}
\label{IV.3.5.2}
Let \(R\) be an equivalence relation in the object \(X\) over \(S\). Let
\(S'\to S\) be a universal effective epimorphism. Suppose that every
\(S\)-morphism whose inverse image on \(S'\) is in \((M)\) is itself in
\((M)\). The following conditions are equivalent:
\begin{enumeratei}
\item \(R\) is \((M)\)-effective in \(X\);
\item \(R'\) is \((M)\)-effective in \(X'\) and the canonical descent datum
on \(X'/R'\) is effective.
\end{enumeratei}
If this is so, the ``descended'' object of \(X'/R'\) is canonically
isomorphic to \(X/R\).
\end{proposition}

The fact that (i)~entails~(ii) is nothing other than the translation into
the language of descent of~3.4.4.%
% typo?: 3.4.4 kept as printed
If one shows the converse, the last assertion of the proposition will be a
consequence of the fact that a universal effective epimorphism is a
descent morphism (2.3), hence that the ``descended object'' is unique (up
to unique isomorphism).

Let us therefore prove (ii)\(\Rightarrow\)(i). Let \(Y'\) be the quotient
\(X'/R'\) and \(Y\) the descended object. As the canonical morphism
\(p'\colon X'\to X'/R'=Y'\) is compatible by construction with the descent
data (its inverse images on \(S''\) coincide with the canonical morphism
\(X''\to X''/R''\)), it arises by inverse image on \(S'\) from an
\(S\)-morphism \(p\colon X\to Y\). As \(p'\) is an element of~\((M)\), it
follows from the hypothesis made on the morphism \(S'\to S\) that \(p\) is
likewise an element of~\((M)\). As \(p'\) is compatible with the
equivalence relation \(R'\), \(p\) is compatible with \(R\), again because
a universal effective epimorphism is a descent morphism. One therefore
has a morphism
\[
R\longrightarrow X\times_{Y}X.
\]
To prove that \(R\) is \((M)\)-effective and that \(Y\) is isomorphic to
\(X/R\), it suffices to prove that this morphism is an isomorphism. Now it
becomes one by extension of the base from \(S\) to \(S'\), for \(R'\) is
effective; it is therefore an isomorphism for the same reason as
previously (2.4).

Let us remark that the hypothesis of the text is verified if one takes
for~\((M)\) the family \((M_{0})\) of universal effective epimorphisms and
if \(\mathcal{C}\) possesses fiber products (1.10).
One deduces from this the

\begin{corollary}{3.5.3}
\label{IV.3.5.3}
Suppose that \(\mathcal{C}\) possesses fiber products (over \(S\) would
suffice). Let \(R\) be an equivalence relation in \(X\) over \(S\) and
\(S'\to S\) a universal effective epimorphism. The following conditions
are equivalent:
\begin{enumeratei}
\item \(R\) is universally effective in \(X\),
\item \(R'\) is universally effective in \(X'\) and the canonical descent
datum on \(X'/R'\) is effective.
\end{enumeratei}
If this is so, the ``descended'' object of \(X'/R'\) is canonically
isomorphic to \(X/R\).
\end{corollary}

\sgasection{4}{Topologies and sheaves}
\label{IV.sec.4}

The notion of sieve, and the presentation of the notion of topology
(4.2.1) adopted here, (more intrinsic and more convenient in many respects
than that by covering families of~[MA]), are due to J.~Giraud~[AS].

\sgasection{4.1}{Sieves}\label{IV.sec.4.1}

\begin{definition}{4.1.1}
\label{IV.4.1.1}
One calls a sieve of the category \(\mathcal{C}\) a subfunctor \(C\) of the
final functor \(\mathbf{e}\colon\mathcal{C}^{\circ}\to\Ens\).
\end{definition}

With every sieve \(C\) of \(\mathcal{C}\) one associates the set \(E(C)\) of
the objects \(X\) of \(\mathcal{C}\) such that \(C(X)\neq\varnothing\), that
is, such that the structural morphism \(X\to\mathbf{e}\) factors
through~\(C\). One therefore has the equivalences
\[
(+)\qquad
\begin{cases}
\,X\in E(C)\iff C(X)=\mathbf{e}(X)=\{\varnothing\}.\\
\,X\notin E(C)\iff C(X)=\varnothing.
\end{cases}
\]
The set \(E=E(C)\) enjoys the following property:

\par\smallskip\noindent
\((++)\)\qquad
If \(X\in E\) and if \(\Hom(Y,X)\neq\varnothing\), then \(Y\in E\).

\noindent
(Let us remark that if one equips the set \(\Ob(\mathcal{C})\) with its
\emph{natural preorder structure},
(\(Y\) \emph{dominating} \(X\) if there exists an arrow from \(Y\) to \(X\)),
the sets \(E\) satisfying \((++)\) are the subsets of \(\Ob(\mathcal{C})\)
which contain every majorant%
{\renewcommand{\thefootnote}{(17)}\nde{: Here, ``majorant'' is taken in
the sense of the preorder relation mentioned above, that is, \(Y\)
majorizes \(X\) if there exists an arrow \(Y\to X\). On the other hand, if
\(X\), \(Y\) are two subobjects of an object \(Z\), one says (\cf~2.4)
that \(Y\) majorizes \(X\) if \(X\subset Y\). To avoid any ambiguity
between these two terminologies, one has in the sequel replaced
``majorant'' by ``dominating'' in the first case, and by ``containing'',
in the second.}}
of one of their elements.)

Conversely, if \(E\) is a subset of \(\Ob(\mathcal{C})\) enjoying the
property~\((++)\), then \(E\) may be written in a unique manner in the form
\(E(C)\) and \(C\) is defined by the formulas~\((+)\). There is therefore a
one-to-one correspondence between the sieves of \(\mathcal{C}\) and the
subsets of \(\Ob(\mathcal{C})\) satisfying the condition~\((++)\). By
\emph{abuse of language}, we shall sometimes say that the set \(E(C)\) is
a sieve of~\(\mathcal{C}\).

{\renewcommand{\thefootnote}{(18)}\nde{: One has spelled out the sentence
which follows.}}%
Let \(C\) and \(C'\) be two sieves of \(\mathcal{C}\); as they are two
subfunctors of the final functor \(\mathbf{e}\), it amounts to the same to
say that \(C\) \emph{dominates} \(C'\) (for the domination relation in
\(\Ob(\widehat{\mathcal{C}})\)), or that \(C\) is a \emph{subfunctor}
of~\(C'\), or again that \(E(C)\subset E(C')\); one will then say that
\(C\) is \emph{finer} than~\(C'\). One sees that one is dealing in this
case with an order structure on the set of sieves of~\(\mathcal{C}\).
Moreover, one has \(E(C)\cap E(C')=E(C\times C')\) and hence the set of
sieves of \(\mathcal{C}\) is \emph{filtered} for the order relation ``being
finer''.

Every subset \(E\) of \(\Ob(\widehat{\mathcal{C}})\), for example a subset
of \(\Ob(\mathcal{C})\), defines a sieve \(C(E)\): the set of the
\(X\in\Ob(\mathcal{C})\) such that \(F(X)\neq\varnothing\) for at least one
\(F\in E\) satisfies the condition~\((++)\) and defines the sieve sought.

This sieve may also be defined as the \emph{image} of the family of
morphisms \(\{F\to\mathbf{e},\ F\in E\}\) in the sense of the following
definition:

\begin{definition}{4.1.2}
\label{IV.4.1.2}
Let \(\{F_{i}\to F\}\) be a family of morphisms of
\(\widehat{\mathcal{C}}\) with the same target~\(F\). One calls the
\emph{image} of this family the subfunctor of \(F\) defined by
\[
S\longmapsto\bigcup_{i}\Im F_{i}(S)\subset F(S).
\]
\end{definition}

\begin{proposition}{4.1.3}
\label{IV.4.1.3}
The formation of the image commutes with extension of the base: in the
preceding notations, let us denote by \(I\) the image of the family
\(\{F_{i}\to F\}\); for every morphism \(G\to F\) of
\(\widehat{\mathcal{C}}\), the image of the family of morphisms
\(\{F_{i}\times_{F}G\to G\}\) is the subfunctor \(I\times_{F}G\) of~\(G\).
\end{proposition}

\begin{definition}{4.1.4.0}
\label{IV.4.1.4.0}
{\renewcommand{\thefootnote}{(19)}\nde{: One has added the
numbering~4.1.4.0, in order to bring this definition into evidence.}}%
Let \(C\) be a sieve of \(\mathcal{C}\). If \(E\) is a subset of
\(\Ob(\mathcal{C})\) such that \(C(E)=C\), one says that \(E\) is a
\emph{base} of~\(C\). Every sieve \(C\) possesses a base, for example the
set \(E(C)\).
\end{definition}

We propose to describe the set \(\Hom(C,F)\), where \(C\) is a sieve of
\(\mathcal{C}\) and \(F\) an object of \(\widehat{\mathcal{C}}\), by means
of a base \(\{S_{i}\}\) of~\(C\). For each pair \((i,j)\), one has a
diagram in \(\widehat{\mathcal{C}}\):
\[
\xymatrix@C=4em@R=2.4em{
S_{i}\times S_{j} \ar[r] \ar[d] & S_{i} \ar[d] \\
S_{j} \ar[r] & \mathbf{e}\,,
}
\]
whence a diagram of sets
\[
\xymatrix@C=3.4em{
\Gamma(F)=\Hom(\mathbf{e},F) \ar[r]^-{\sigma}
  & \displaystyle\prod_{i}\Hom(S_{i},F)
      \ar@<0.65ex>[r]^-{\tau_{1}}
      \ar@<-0.65ex>[r]_-{\tau_{2}}
  & \displaystyle\prod_{i,j}\Hom(S_{i}\times S_{j},F)
}
\]
such that \(\tau_{1}\sigma=\tau_{2}\sigma\). One therefore has a morphism
\[
\Hom(\mathbf{e},F)\longrightarrow
\Ker\Biggl(\prod_{i}\Hom(S_{i},F)
  \Rightarrow
  \prod_{i,j}\Hom(S_{i}\times S_{j},F)\Biggr).
\]
One verifies immediately:

\begin{proposition}{4.1.4}
\label{IV.4.1.4}
One has an isomorphism, functorial in~\(F\),
\[
\Hom(C,F)\isomto
\Ker\Biggl(\prod_{i}\Hom(S_{i},F)
  \Rightarrow
  \prod_{i,j}\Hom(S_{i}\times S_{j},F)\Biggr),
\]
such that the diagram
\[
\xymatrix@C=2.2em@R=2.8em{
\Hom(\mathbf{e},F) \ar[r]
  & \Ker\Bigl(\prod_{i}\Hom(S_{i},F)
      \Rightarrow
      \prod_{i,j}\Hom(S_{i}\times S_{j},F)\Bigr)
\\
\Hom(\mathbf{e},F) \ar[r]
  & \Hom(C,F) \ar[u]_{\wr}
\,,}
\]
where the last line is induced by the canonical morphism
\(C\to\mathbf{e}\), is commutative.
\end{proposition}

\begin{corollary}{4.1.5}
\label{IV.4.1.5}
Suppose that the fiber products \(S_{i}\times S_{j}\) are representable,
for example that the \(S_{i}\) be quarrable. One then has, for every
\(F\in\Ob(\widehat{\mathcal{C}})\), an isomorphism
\[
\Hom(C,F)\isomto
\Ker\Biggl[\prod_{i}F(S_{i})
  \Rightarrow
  \prod_{i,j}F(S_{i}\times S_{j})\Biggr].
\]
\end{corollary}

\begin{remark}{4.1.6}
\label{IV.4.1.6}
Let \(R\) be a sieve of \(\mathcal{C}\); let us denote by \(\mathcal{R}\)
the full subcategory of \(\mathcal{C}\) whose set of objects is \(E(R)\)
and by
\[
i_{R}\colon\mathcal{R}\longrightarrow\mathcal{C}
\]
the inclusion functor. One has an isomorphism, functorial in
\(F\in\Ob(\widehat{\mathcal{C}})\),
\[
\Hom(R,F)\isomto\Gamma(F\circ i_{R})
\]
such that the diagram
\[
\xymatrix@C=4.2em@R=2.6em{
\Hom(\mathbf{e},F) \ar[r] \ar[d]_{\wr}
  & \Hom(R,F) \ar[d]^{\wr} \\
\Gamma(F) \ar[r]
  & \Gamma(F\circ i_{R})
\,,}
\]
where the second line is induced by the functor \(i_{R}\), is commutative.
\end{remark}

\begin{definition}{4.1.7}
\label{IV.4.1.7}
Let \(\mathcal{C}\) be a category. One calls a sieve of the object \(S\) of
\(\mathcal{C}\) a sieve of the category \(\mathcal{C}_{/S}\).
\end{definition}

A sieve of \(S\) is therefore a sub-\(\widehat{\mathcal{C}}\)-object
of~\(S\). There corresponds to it canonically a subset of
\(\Ob(\mathcal{C}_{/S})\) containing the source of every arrow whose target
it contains. By \emph{abuse of language}, such a set will also be called a
sieve of~\(S\).

\sgasection{4.2}{Topologies: definitions}\label{IV.sec.4.2}

\begin{definition}{4.2.1}
\label{IV.4.2.1}
Let \(\mathcal{C}\) be a category. One calls a topology on \(\mathcal{C}\)
the giving, for each \(S\) of \(\mathcal{C}\), of a set \(J(S)\) of sieves
of~\(S\), called \emph{covering} sieves or \emph{refinements} of~\(S\),
this giving satisfying the following axioms:
\begin{enumerate}[label=\textup{(T~\arabic*)},leftmargin=2.8em,itemsep=0.45em,topsep=0.45em]
\item For every refinement \(R\) of \(S\) and every morphism \(T\to S\), the
sieve \(R\times_{S}T\) of \(T\) is \emph{covering} (``stability under
\emph{base change}'').
\item If \(R\), \(C\) are two sieves of \(S\), if \(R\) is covering and if
for every \(T\in\Ob(\mathcal{C})\) and every morphism \(T\to R\), the sieve
\(C\times_{S}T\) of \(T\) is covering, then \(C\) is a refinement
of~\(S\).%
{\renewcommand{\thefootnote}{(20)}\nde{: that is: if \(C\) is covering
``locally with respect to the covering sieve \(R\)'', then \(C\) is
covering.}}
\item If \(C\supset R\) are two sieves of \(S\) and if \(R\) is covering,
then \(C\) is covering.
\item For every \(S\), \(S\) is a refinement of \(S\).
\end{enumerate}
\end{definition}

One may reformulate these axioms in the following manner. Suppose given a
topology \(S\mapsto J(S)\) on \(\mathcal{C}\) and, for every object \(F\) of
\(\widehat{\mathcal{C}}\), let us denote by \(J(F)\) the set of
subfunctors \(R\) of \(F\) such that for every morphism \(T\to F\) of
\(\widehat{\mathcal{C}}\), where \(T\) is \emph{representable},
\(R\times_{F}T\), which is a sieve of \(T\), is \emph{covering}. By virtue
of~(T~1), this notation is indeed compatible with the preceding one. One
will likewise say that \(R\in J(F)\) is a \emph{refinement} of~\(F\). One
verifies immediately that the preceding axioms entail the following
properties:

\par\smallskip\noindent
(\(T'\)~0) If \(F\supset G\) are two objects of \(\widehat{\mathcal{C}}\),
and if for every \(S\in\Ob(\mathcal{C})\) and every morphism \(S\to F\),
\(G\times_{F}S\in J(S)\), then \(G\in J(F)\).

\par\smallskip\noindent
(\(T'\)~1) If \(G\in J(F)\), and if \(H\to F\) is a morphism of
\(\widehat{\mathcal{C}}\), then \(G\times_{F}H\in J(H)\).

\par\smallskip\noindent
(\(T'\)~2) If \(F\supset G\supset H\) are three objects of
\(\widehat{\mathcal{C}}\), if \(G\in J(F)\) and \(H\in J(G)\), then
\(H\in J(F)\).

\par\smallskip\noindent
(\(T'\)~3) If \(F\supset G\supset H\) are three objects of
\(\widehat{\mathcal{C}}\) and if \(H\in J(F)\), then \(G\in J(F)\).

\par\smallskip\noindent
(\(T'\)~4) For every \(F\in\Ob(\widehat{\mathcal{C}})\), \(F\in J(F)\).

Conversely, if one gives oneself for every
\(F\in\Ob(\widehat{\mathcal{C}})\) a set \(J(F)\) of subobjects of \(F\)
satisfying the properties (\(T'\)~0) to (\(T'\)~4), the map
\(S\mapsto J(S)\) defines a topology on \(\mathcal{C}\) and the two
preceding constructions are inverse to one another.

From (\(T'\)~1), (\(T'\)~2) and (\(T'\)~3)%
{\renewcommand{\thefootnote}{(21)}\nde{: (\(T'\)~1) and (\(T'\)~2)
suffice: \(G\cap H=G\times_{F}H\) belongs to \(J(H)\), by (\(T'\)~1), hence
to \(J(F)\), by (\(T'\)~2).}}
there results the following property:

\par\smallskip\noindent
(\(T'\)~5) If \(G\) and \(H\) are two subobjects of \(F\) and if
\(G,\,H\in J(F)\), then \(G\cap H\in J(F)\).

\noindent
The set \(J(F)\), ordered by the relation \(\supset\), is therefore
\emph{filtered}; this remark will be of use to us later.

\numpara{4.2.2}\label{IV.4.2.2}
One will say that the topology defined by \(J\) is finer than the topology
defined by \(J'\) if for every \(S\in\Ob(\mathcal{C})\),
\(J(S)\supset J'(S)\) (it amounts to the same to say that for every
\(F\in\Ob(\widehat{\mathcal{C}})\), \(J(F)\supset J'(F)\)).

Every set of topologies on \(\mathcal{C}\) possesses a greatest
\emph{lower} bound: let \(I\) be a set of indices,
and for each \(i\in I\), let \(S\mapsto J_{i}(S)\) be a topology on
\(\mathcal{C}\). Set \(J(S)=\bigcap_{i\in I}J_{i}(S)\); it is immediate
that one has thus defined a topology on \(\mathcal{C}\), and that it is
indeed the greatest \emph{lower} bound of the given set.

In particular, let us give ourselves, for each \(S\in\Ob(\mathcal{C})\), a
set \(E(S)\) of sieves of~\(S\). One calls the topology \emph{generated} by
these sets the \emph{least fine} topology for which the elements of
\(E(S)\) are refinements of \(S\) for every~\(S\).

\begin{definition}{4.2.3}
\label{IV.4.2.3}
Let \(\{F_{i}\to F\}\) be a family of morphisms of
\(\widehat{\mathcal{C}}\). Let \(G\subset F\) be the image (4.1.2) of this
family. The family is said to be \emph{covering} if \(G\in J(F)\). A
morphism is said to be \emph{covering} if the family reduced to this
morphism is covering.
\end{definition}

This definition applies in particular to an inclusion: a sieve \(C\) of
\(S\) is covering if and only if the canonical morphism \(C\to S\) is
covering.

The axioms (\(T'\)~0) to (\(T'\)~5) entail for covering families the
following properties:
